%\usepackage[letterpaper,showframe]{geometry}	
% usuwanie nagłówków/stopek z\,dokumentu
\usepackage[indent=0pt]{parskip} % ustawianie paragrafów
%\setstretch{1.5} % półtorej linii odległości między wierszami
\usepackage{titlesec} % ustawianie nagłówków
%\usepackage{fullpage} % redukcja marginesów
\usepackage{setspace} % odstępy między liniami
\usepackage{enumitem}% listy numerowane
\usepackage{soul}% dodaje \ul
\usepackage{ulem}% dodaje \uline



% -----------------------------
\renewcommand*{\thepage}{\footnotesize{\arabic{page}}}% mniejsza czcionka numeru strony, dodanie czegokolwiek przed numerem: {A--\arabic{page}}
\newcommand{\margintext}[1]{\textit{\small #1}}
% przesunięcia poziome
\newcommand{\smallwidth}{\hskip 1em}
\newcommand{\medwidth}{\hskip 3pt}
\newcommand{\largewidth}{\hskip 12pt}
% przesunięcia pionowe
\newcommand{\linespace}{\vskip 7pt}
\newcommand{\paragraphspace}{\vskip 8pt}
\newcommand{\segmentspace}{\vskip 18pt}
\newcommand{\sectionspace}{\vskip 28pt}
%------------------------------------------

% --- GREGORIO ---	
\usepackage{gregoriotex}
\gresetlinecolor{gregoriocolor}% Make staff lines red; remove for black:
%\gresetgregoriofont{granapadano} %{gregorio}
\gresetheadercapture{annotation}{greannotation}{string}% odczytuje adnotacje nad inicjałem jako ciąg, co pozwala na jego renderowanie
\gresetheadercapture{commentary}{grecommentary}{string}% odczytuje komentarz jako ciąg, co pozwala na jego renderowanie
\grechangedim{commentaryraise}{0.3cm}{scalable} % podnosi komentarz wystarczająco ponad górną linię
\grechangestyle{annotation}{\color{black}\small}% format adnotacji
\grechangestyle{commentary}{\color{gregoriocolor}\small}% format komentarza
\grechangestyle{abovelinestext}{\color{gregoriocolor}\small}% format napisu nad liniami
\grechangestyle{initial}{\fontsize{36}{36}\selectfont} % zmniejsza wielkość inicjału do bardziej proporcjonalnego rozmiaru (=36pt)
\grechangedim{intersyllablespacenotes}{0.05cm}{scalable}% Minimalna odległość między nutami różnych sylab.
% Set the space around the initial:
% See http://gregorio-project.github.io/gregoriotex/details.html for more details and options
\grechangedim{beforeinitialshift}{0.5mm}{scalable}
\grechangedim{afterinitialshift}{1.0mm}{scalable}
\grechangestaffsize{25}%rozmiar wyniku GregorioTEX

% umieszczenie zawartości składu ich unikalnego argumentu po lewej i prawej stronie każdej linii
\let\grelocalleftbox\localleftbox
\let\grelocalrightbox\localrightbox

%\newcommand{\separator}{\greseparator{3}{40}}
%\newfontface\vectis{VectisMonumentalW01-Regular}
%\def\greinitialformat#1{\fontsize{20}{20}\vectis #1}
% ------------------

\usepackage[hyperindex=true]{hyperref}

%---- usun konflikt między paracol i luacolor ----
\wstawplik{conparlua}
%-----------------------------------------------

% --- SPIS TREŚCI -------
%\usepackage{imakeidx}
%\makeindex[title=Alfabetyczny spis program=texindy,options=-M styl.xdy]
%\usepackage{filecontents}
%\begin{filecontents*}{styl.xdy}
%;;; xindy style file
%(markup-locclass-list :open "\dotfill" :sep "")
%(define-attributes (( "textbf" "default" )) )
%(markup-locref   :attr  "textbf"     :open "\textbf{\hyperpage{" :close "}}")
%(markup-locref   :attr  "textit"     :open "\textit{\hyperpage{" :close "}}")
%(markup-locref   :attr  "textttt"     :open "\textttt{\hyperpage{" :close "}}")
%(markup-locref   :attr  "texttsc"     :open "\texttsc{\hyperpage{" :close "}}")
%(markup-locref   :attr  "default"     :open "\hyperpage{" :close "}")
%\end{filecontents*}
% -------------------------

% ---  Konfiguracja kolumn ---
\usepackage{paracol}
\setcolumnwidth{.49\textwidth, .49\textwidth}
\setlength{\columnseprule}{0.4pt}
\setlength{\columnsep}{.02\textwidth}
\colseprulecolor{red}% wymagany pakiet xcolor
\footnotelayout{m}
% ----------------------------

\makeatletter 
\patchcmd\pcol@hfil
 {\copy\pcol@tempboxa}
 {\luacolorresetboxcolor\pcol@tempboxa\copy\pcol@tempboxa}{}{\fail}
\makeatother
% ----------------------------

\footnotelayout{m}

%% --- UKŁAD DWUKOLUMNOWY Z TŁUMACZENIEM ---
%\usepackage{tcolorbox}
%\newtcolorbox{nutybox}{size=minimal, colback=white, colframe=red, arc=0mm, width=\linewidth, halign=left, valign=center}
%\newtcolorbox{tlumbox}{size=minimal, colback=white, colframe=red, arc=0mm, width=\linewidth, halign=left, valign=top, fontupper=\small\itshape} 
%\newcommand{\nuty}[1]{\begin{nutybox}#1 \end{nutybox}}
%\newcommand{\tlum}[1]{\begin{tlumbox}\tlumacz{#1}\end{tlumbox}}
%% ------------------------------------------------------

\usepackage{Tabbing}
\usepackage{longtable}
\usepackage[savepos]{zref}

% --- CZCIONKI DOKUMENTU ---
%\usepackage{fontspec}
%\defaultfontfeatures{Ligatures=TeX}
\usepackage{polyglossia}
%\usepackage{ecclesiastic}% This package works only with pdfLaTeX 
\usepackage[autostyle=true]{csquotes}%generuje odpowiednie cudzysłowy dla danego języka
\setdefaultlanguage{polish}
\setotherlanguage[variant=ecclesiastic, hyphenation=liturgical]{latin}

%\usepackage[latin,polish]{babel}
%\usepackage{times}
%\usepackage[math]{anttor}
%\usepackage{tgtermes}
%\usepackage{mathptmx}
\usepackage{ebgaramond}
\usepackage[osf,p]{libertine}


% definiowanie fontów w  LuaLATEX':
\setmainfont{EB Garamond}
%\setsansfont{CMU Sans Serif}
%\setmonofont{CMU Typewriter Text}
% ----
%\setmainfont{TeX Gyre Termes}
%\setmainfont{TeX Gyre Heros}
%\setmainfont{Calluna Regular}

%\setmainfont{Bodoni Moda}
%\newfontfamily{\sekfont}{Times New Roman}
\newfontfamily{\sekfont}{Libertinus Serif}
\newcommand\sekcjafont{\fontsize{16}{16}\selectfont \sekfont}

%----- Strona tytułowa -------------~*
\newfontfamily{\tytfont}{JSL Ancient}
\newcommand\tytulfont[1]{\fontsize{36}{30}\selectfont \tytfont#1\par}
\newcommand\podtytulfont[1]{\fontsize{18}{10}\selectfont \tytfont#1\par}

%\setmainfont[<font features>]{<name of font>} % sets the roman font
%\setsansfont[<font features>]{<name of font>} % sets the sans font
%\setmonofont[<font features>]{<name of font>} % sets the monospace font

% --- KONIEC CZCIONKI DOKUMENTU ---

% --- GRAFIKA I KOLORY ---
\usepackage{xcolor} % kolory
\usepackage{graphicx} % umożliwia import obrazów
\graphicspath{{img/}} % katalog obrazów
\usepackage{pgfornament}% ornamenty graficzne
% --- KONIEC GRAFIKA I KOLORY -----------

% --- KOMENDY WŁASNE ---
\newcommand{\red}[1]{\textcolor{red}{#1}}
\newcommand{\black}[1]{\textcolor{black}{#1}}
% --------------------------------
%\newcommand{\firstchar}[1]{\directlua{tex.sprint(unicode.utf8.sub('#1',1,1))}}%makro pobiera pierwszy znak ze stringu
\newcommand{\tlumacz}[1]{\small\textit{#1}}
\newcommand{\inicjal}[2]{\textcolor{red}{\Large\textbf{#1}}{#2}}
\newcommand{\tlumaczenie}[3]{\selectlanguage{polish}\textcolor{red}{\Large\textbf{#1}}{#2}{#3}}%
%--------------------------------
% ----------- RUBRYKA -------------------
\newcommand{\rubryka}[1]{
	\begin{spacing}{0.8}
	\begin{center}
		\parbox{0.9\textwidth}{\centering\small\textit{\red{#1}}}
	\end{center}
	\end{spacing}
	\linespace}
%-----------STYLE NAGŁÓWKÓW ----------------
%
% ----------- ROZDZIAŁ --------------------
\titlespacing*{\chapter}{0pt}{-40pt}{0pt}	
\titleformat{\chapter}[hang]{\filcenter\scshape\bfseries\Large}{}{0pt}{} % format nagłówków
\newcommand{\rozdzial}[2]{
		\chapter{#1}
		\vspace{-1em}
		\rubryka{#2}}

% ----------- SEKCJA --------------------
\titlespacing*{\section}{0pt}{0pt}{0pt}	
\titleformat{\section}[hang]{\filcenter\scshape\normalfont}{}{0pt}{} % format nagłówków
\newcommand{\sekcja}[2]{
	\begin{spacing}{0.6}
	\begin{center}
		\segmentspace
		\section*{#1}
		\begin{minipage}[c]{0.88\linewidth}
			\rubryka{#2}
			%\centering\red{\small\textit{#2}}
		\end{minipage}
	\end{center}	
	\end{spacing}
	\linespace}
	
% ----------- SUBSEKCJA --------------------	
\titlespacing*{\subsection}{0pt}{0pt}{0pt}	
\titleformat{\subsection}[hang]{\normalfont\scshape}{}{0pt}{} % format podnagłówków
\newcommand{\subsekcja}[1]{
	\segmentspace
	\subsection*{\red{#1}}		
	\segmentspace}
%------------ Tytuł -----------------------
\newcommand{\tytul}[1]{
	\red{\large\scshape{#1}\linespace}}
%------------ Podytuł -----------------------
\newcommand{\podtytul}[1]{
	\red{\normalfont\scshape{#1}\linespace}}
%------------DEKOR ---------------------
\newcommand{\dekor}[2]{
	\begin{center}
		\vspace*{\fill}
		\includegraphics[width=#1\textwidth]{#2}
		\vspace*{\fill}
	\end{center}}
%------- Zarządzanie wersami w psalmach z tlumaczeniem -------------------------
\newcounter{licznik}
\newcommand{\licznikZero}{\setcounter{licznik}{0}}% zeruj licznik
\newcommand{\licznikUp}{\refstepcounter{licznik}}% zwiększ licznik 
\licznikZero
\newcommand{\BeginParallel}{\begin{spacing}{0.8}\begin{paracol}{2}\licznikZero}
\newcommand{\lat}[1]{\selectlanguage{latin}\ensurevspace{2\baselineskip}\begin{minipage}[t]{\linewidth}\ifnum \value{licznik}>0 \licznikUp \thelicznik .~#1 \else \licznikUp#1 \fi \end{minipage}\switchcolumn}
\newcommand{\latN}[1]{\selectlanguage{latin}\ensurevspace{2\baselineskip}\begin{minipage}[t]{\linewidth}#1\end{minipage}\switchcolumn}% kolumna bez numeracji wersów
\newcommand{\pol}[1]{\selectlanguage{polish}\begin{minipage}[t]{\linewidth}#1\end{minipage}\switchcolumn*}
\newcommand{\EndParallel}{\end{paracol}\end{spacing}}

% --- ZNAKI SPECJALNE ---
% 
%\newcommand †
% --------------------------
